% ================================================================
%  SESIÓN 17 — Seleccionar la información de la guía
%  Almería en Cifras y Formas · Matemáticas 1.º ESO · IES Al-Ándalus
%  Compilar: pdflatex sesion17.tex
% ================================================================
\documentclass[aspectratio=169, 12pt, spanish]{beamer}
\usepackage{almeria-beamer}

\renewcommand{\SessionNumber}{17}
\renewcommand{\SessionTitle}{Seleccionar la información\\de la guía}
\renewcommand{\SessionName}{Sesión 17}
\renewcommand{\SessionObjective}{Aplicar criterios matemáticos y comunicativos para seleccionar, verificar y organizar la información que formará parte de la guía turística matemática de Almería}
\renewcommand{\BlockName}{Bloque 4: Producto final}

\begin{document}

% ── PORTADA ────────────────────────────────────────────────────────
\begin{frame}[plain]
  \titlepage
\end{frame}

% ── ÍNDICE ─────────────────────────────────────────────────────────
\begin{frame}{Contenidos de hoy}
  \begin{columns}[t]
    \begin{column}{0.48\textwidth}
      \begin{enumerate}
        \item Contexto: ¿qué hemos recogido?
        \item Criterios de selección
        \item Lista de verificación de datos
        \item Problemas de calidad habituales
        \item Árbol de decisión
      \end{enumerate}
    \end{column}
    \begin{column}{0.48\textwidth}
      \begin{enumerate}\setcounter{enumi}{5}
        \item Categorías de la guía
        \item Evaluar datos concretos
        \item Ejercicios ANALIZAR / EVALUAR / CREAR
        \item Evidencia del día
      \end{enumerate}
    \end{column}
  \end{columns}
  \vspace{0.4cm}
  \begin{center}
    \bloomtag{Analizar}\bloomtag{Evaluar}\bloomtag{Crear}
  \end{center}
\end{frame}

% ── SECCIÓN 1 ──────────────────────────────────────────────────────
\seccionframe[BAnalizar]{De la exploración a la selección}

% ── FRAME 3: Contexto ──────────────────────────────────────────────
\begin{frame}{Ya tenemos datos: ¿cuáles elegimos?}
  \begin{columns}[c]
    \begin{column}{0.52\textwidth}
      \begin{teoriabox}
        En las \textbf{16 sesiones anteriores} hemos recogido datos matemáticos sobre Almería:
        \begin{itemize}
          \item Coordenadas de monumentos
          \item Áreas, perímetros y alturas
          \item Tablas estadísticas y gráficas
          \item Rutas y distancias
          \item Relaciones de proporcionalidad
        \end{itemize}
        \vspace{0.2cm}
        Ahora toca \textbf{seleccionar lo mejor} para la guía.
      \end{teoriabox}
    \end{column}
    \begin{column}{0.45\textwidth}
      \begin{center}
      \begin{tikzpicture}[scale=0.9]
        % Triángulo de selección
        \node[draw=AlmAzul, fill=AlmAzulC!15, rounded corners=6pt,
              minimum width=3.5cm, minimum height=0.65cm,
              font=\scriptsize\bfseries, color=AlmAzul] (a) at (0, 2.2)
          {Muchos datos (sesiones 1-16)};
        \node[draw=AlmAmbar!80!black, fill=AlmAmbar!15, rounded corners=6pt,
              minimum width=2.6cm, minimum height=0.65cm,
              font=\scriptsize\bfseries, color=AlmAzul] (b) at (0, 1.1)
          {Datos candidatos};
        \node[draw=AlmVerde, fill=AlmVerde!15, rounded corners=6pt,
              minimum width=2.0cm, minimum height=0.65cm,
              font=\scriptsize\bfseries, color=AlmVerde] (c) at (0, 0)
          {Datos seleccionados};
        \draw[->, thick, AlmAzulM] (a) -- (b);
        \draw[->, thick, AlmAmbar!70!black] (b) -- (c);
        \node[right=0.1cm of a, font=\tiny, color=AlmGris] {100\%};
        \node[right=0.1cm of b, font=\tiny, color=AlmGris] {50\%};
        \node[right=0.1cm of c, font=\tiny, color=AlmGris] {25\%};
      \end{tikzpicture}
      \end{center}
      \begin{notabox}[Principio]
        \textbf{Calidad} sobre cantidad.
      \end{notabox}
    \end{column}
  \end{columns}
\end{frame}

% ── FRAME 4: Criterios de selección ───────────────────────────────
\begin{frame}{Los cinco criterios de selección}
  \begin{columns}[t]
    \begin{column}{0.48\textwidth}
      \begin{block}{\faCheckDouble\ R — Relevancia}
        ¿Es este dato \textbf{matemáticamente interesante}? ¿Aporta algo al turista?
      \end{block}
      \vspace{0.2cm}
      \begin{block}{\faSearch\ E — Exactitud}
        ¿Es el dato \textbf{preciso y verificable}? ¿Tiene fuente?
      \end{block>}
      \vspace{0.2cm}
      \begin{block}{\faComments\ C — Claridad}
        ¿Se puede explicar a un turista \textbf{sin formación matemática}?
      \end{block}
    \end{column}
    \begin{column}{0.48\textwidth}
      \begin{block}{\faMapMarkerAlt\ R — Representatividad}
        ¿Representa bien a Almería? ¿Es \textbf{característico} de la ciudad?
      \end{block}
      \vspace{0.2cm}
      \begin{block}{\faStar\ O — Originalidad}
        ¿Es \textbf{sorprendente o único}? ¿Algo que el turista no esperaría?
      \end{block}
      \vspace{0.3cm}
      \begin{formulabox}[BAnalizar]
        \textbf{R\,·\,E\,·\,C\,·\,R\,·\,O}
      \end{formulabox}
    \end{column}
  \end{columns}
\end{frame}

% ── FRAME 5: Lista de verificación ────────────────────────────────
\begin{frame}{Lista de verificación para cada dato}
  \begin{columns}[c]
    \begin{column}{0.55\textwidth}
      \renewcommand{\arraystretch}{1.4}
      \begin{tabular}{|c|p{6.5cm}|}
        \hline
        \rowcolor{AlmAzul!15}
        \textbf{Check} & \textbf{Pregunta} \\
        \hline
        $\square$ & Fuente verificada (no inventada) \\
        $\square$ & Unidades correctas (m, no m² para distancias) \\
        $\square$ & Cálculos comprobados dos veces \\
        $\square$ & Visualizable (mapa, tabla o gráfica) \\
        $\square$ & Interesante para un turista \\
        \hline
      \end{tabular}
    \end{column}
    \begin{column}{0.42\textwidth}
      \begin{notabox}[Regla de los 5/5]
        Un dato de \textbf{calidad} supera \textbf{todos} los checks.\\[6pt]
        Un dato con 3 o menos checks debe \textbf{revisarse o descartarse}.
      \end{notabox}
      \vspace{0.2cm}
      \begin{notabox}[Unidades]
        Distancia $\to$ m, km\\
        Área $\to$ m², ha, km²\\
        Tiempo $\to$ min, h\\
        Temperatura $\to$ °C
      \end{notabox}
    \end{column}
  \end{columns}
\end{frame}

% ── FRAME 6: Problemas de calidad ─────────────────────────────────
\begin{frame}{Errores habituales en la calidad de los datos}
  \begin{columns}[t]
    \begin{column}{0.48\textwidth}
      \begin{errorbox}
        \textbf{Tipos de datos problemáticos:}
        \begin{itemize}
          \item[\incorrecto] \textbf{Obsoletos:} «La población de Almería en 2000 era...»
          \item[\incorrecto] \textbf{Sin fuente:} «Dicen que la Alcazaba tiene 43.000 m²»
          \item[\incorrecto] \textbf{Sin unidades:} «La torre mide 55»
          \item[\incorrecto] \textbf{Cálculo incorrecto:} perímetro calculado como área
        \end{itemize}
      \end{errorbox}
    \end{column}
    \begin{column}{0.48\textwidth}
      \begin{exampleblock}{\correcto\ Datos de calidad:}
        \begin{itemize}
          \item «La Alcazaba ocupa \textbf{4,3 ha} = 43.000 m²\\ (Fuente: Junta de Andalucía, 2023)»
          \item «El Cabo de Gata tiene \textbf{32 m} de altura (IGN, 2022)»
          \item «Almería recibió \textbf{1,2 millones} de turistas en 2023 (INE)»
        \end{itemize}
      \end{exampleblock}
    \end{column}
  \end{columns}
\end{frame}

% ── FRAME 7: Árbol de decisión TikZ ───────────────────────────────
\begin{frame}{Árbol de decisión: ¿incluyo este dato?}
  \begin{center}
  \begin{tikzpicture}[
    decision/.style={diamond, draw=AlmAzulM, fill=AlmAzulC!20,
                     text width=2.6cm, align=center, font=\tiny\bfseries,
                     inner sep=1pt, aspect=1.8},
    result/.style={rounded corners=4pt, draw=AlmVerde, fill=AlmVerde!15,
                   font=\tiny\bfseries, text width=2cm, align=center, inner sep=3pt},
    reject/.style={rounded corners=4pt, draw=AlmTerra, fill=AlmTerra!10,
                   font=\tiny\bfseries, text width=1.8cm, align=center, inner sep=3pt},
    arr/.style={->, >=Stealth, thick, AlmAzulM},
    scale=0.9, every node/.style={font=\tiny}
  ]
    \node[decision] (Q1) at (0,0) {¿Tiene fuente verificable?};
    \node[reject] (R1) at (-3.2,-1.5) {Descartar o buscar fuente};
    \node[decision] (Q2) at (2.0,-1.5) {¿Tiene unidades correctas?};
    \node[reject] (R2) at (4.8,-1.5) {Corregir unidades};
    \node[decision] (Q3) at (2.0,-3.0) {¿Es relevante para el turista?};
    \node[reject] (R3) at (5.0,-3.0) {Descartar};
    \node[result] (OK) at (2.0,-4.5) {\correcto\ Incluir en la guía};

    \draw[arr] (Q1) -- node[above,sloped]{\tiny No} (R1);
    \draw[arr] (Q1) -- node[above]{\tiny Sí} (Q2);
    \draw[arr] (Q2) -- node[above]{\tiny No} (R2);
    \draw[arr] (Q2) -- node[right]{\tiny Sí} (Q3);
    \draw[arr] (Q3) -- node[above]{\tiny No} (R3);
    \draw[arr] (Q3) -- node[right]{\tiny Sí} (OK);
  \end{tikzpicture}
  \end{center}
\end{frame}

% ── FRAME 8: Categorías de la guía ────────────────────────────────
\begin{frame}{Categorías de la guía turística matemática}
  \begin{columns}[t]
    \begin{column}{0.48\textwidth}
      \begin{definicionbox}{Estructura de la guía}
        \begin{enumerate}
          \item \textbf{Mapa de la ciudad} — coordenadas, escala, puntos clave
          \item \textbf{Monumentos principales} — áreas, alturas, fechas, visitantes
          \item \textbf{Rutas turísticas} — distancias, tiempos, dificultad
        \end{enumerate}
      \end{definicionbox}
    \end{column}
    \begin{column}{0.48\textwidth}
      \begin{definicionbox}{Estructura (cont.)}
        \begin{enumerate}\setcounter{enumi}{3}
          \item \textbf{Datos curiosos} — estadísticas sorprendentes de Almería
          \item \textbf{Gráficas estadísticas} — turismo, temperatura, población
          \item \textbf{Glosario matemático} — términos usados en la guía
        \end{enumerate}
      \end{definicionbox}
    \end{column}
  \end{columns}
  \vspace{0.35cm}
  \begin{center}
  \begin{tikzpicture}
    \foreach \i/\cat/\col in {
      1/Mapa/AlmAzulM,
      2/Monumentos/AlmTerra,
      3/Rutas/AlmVerde,
      4/Curiosidades/AlmAmbar!80!black,
      5/Gráficas/BAnalizar,
      6/Glosario/BCrear}{
      \node[draw=\col, fill=\col!12, rounded corners=3pt,
            font=\tiny\bfseries\color{\col}, inner sep=4pt,
            minimum width=1.8cm, minimum height=0.55cm]
        at ({(\i-1)*2.2},0) {\cat};
    }
  \end{tikzpicture}
  \end{center}
\end{frame>}

% ── FRAME 9: Evaluar datos concretos ──────────────────────────────
\begin{frame}{Evaluando datos concretos}
  \begin{columns}[t]
    \begin{column}{0.48\textwidth}
      \begin{block}{Dato A}
        «La Alcazaba ocupa aproximadamente 43.000 m²»
        \vspace{0.2cm}

        \begin{tabular}{lc}
          Fuente verificable & $\square$ \\
          Unidades correctas & \correcto \\
          Relevante & \correcto \\
          Verificable & \correcto \\
        \end{tabular}
        \vspace{0.2cm}
        \textcolor{AlmAmbar!80!black}{\textbf{Candidato: buscar fuente}}
      \end{block}
    \end{column}
    \begin{column}{0.48\textwidth}
      \begin{alertblock}{Dato B}
        «La Alcazaba tiene 100 torres»
        \vspace{0.2cm}

        \begin{tabular}{lc}
          Fuente verificable & \incorrecto \\
          Unidades correctas & \correcto \\
          Relevante & \correcto \\
          Plausible & \incorrecto \\
        \end{tabular}
        \vspace{0.2cm}
        \textcolor{AlmTerra}{\textbf{Rechazado: valor sospechoso}}
      \end{alertblock}
    \end{column}
  \end{columns}
\end{frame}

% ── FRAME 10: Ejercicio ANALIZAR ──────────────────────────────────
\begin{frame}{Ejercicio — \bloomtag{Analizar}}
  \begin{analizarbox}
    \textbf{8 candidatos para la guía.} Aplica la lista de verificación a cada uno y decide cuáles de los 5 incluirías:

    \vspace{0.2cm}
    \begin{columns}[t]
      \begin{column}{0.48\textwidth}
        \begin{enumerate}
          \item «El Paseo de Almería mide 1,2 km de largo» (sin fuente)
          \item «El faro del Cabo de Gata: 32 m» (IGN 2022)
          \item «Almería tiene muchos días de sol» (sin número)
          \item «La temperatura media en julio es 27°C» (AEMET 2023)
        \end{enumerate}
      \end{column>}
      \begin{column}{0.48\textwidth}
        \begin{enumerate}\setcounter{enumi}{4}
          \item «El Cable Inglés mide 1.100 m» (Ayto. 2021)
          \item «La Alcazaba es muy grande» (sin dato)
          \item «En 2023 llegaron 1.200.000 turistas» (INE)
          \item «El estadio Juegos Mediterráneos: 15.000 localidades» (sin año)
        \end{enumerate}
      \end{column}
    \end{columns}

    \vspace{0.2cm}
    \textbf{¿Cuáles 5 incluyes? ¿Por qué?} Usa la rúbrica R-E-C-R-O.
  \end{analizarbox}
\end{frame>}

% ── FRAME 11: Ejercicio EVALUAR ───────────────────────────────────
\begin{frame}{Ejercicio — \bloomtag{Evaluar}}
  \begin{evaluarbox}
    \textbf{Tu equipo tiene los siguientes materiales:}
    \begin{columns}[t]
      \begin{column}{0.48\textwidth}
        \begin{itemize}
          \item 3 mapas con coordenadas
          \item 2 gráficas estadísticas
          \item 5 listas de coordenadas
          \item 1 tabla de cálculos verificados
        \end{itemize}
      \end{column}
      \begin{column}{0.48\textwidth}
        Decide la \textbf{estructura óptima} de tu capítulo de guía:
        \begin{enumerate}
          \item ¿Qué materiales incluyes?
          \item ¿Cuáles descartas? ¿Por qué?
          \item ¿Qué sección del capítulo encabeza cada material?
          \item ¿Falta algo? ¿Qué añadirías?
        \end{enumerate}
      \end{column}
    \end{columns}
    \vspace{0.2cm}
    Justifica \textbf{cada} decisión con criterios matemáticos.
    \dificultad{3}
  \end{evaluarbox}
\end{frame>}

% ── FRAME 12: Ejercicio CREAR ─────────────────────────────────────
\begin{frame}{Ejercicio — \bloomtag{Crear}}
  \begin{crearbox}
    \textbf{Diseña la portada de tu guía turística matemática:}

    \vspace{0.3cm}
    \begin{columns}[t]
      \begin{column}{0.55\textwidth}
        La portada debe incluir:
        \begin{enumerate}
          \item \textbf{Título} de la guía
          \item \textbf{Subtítulo} que describa el contenido matemático
          \item \textbf{Una estadística clave} sobre Almería (bien elegida y verificada)
          \item \textbf{Miniatura de mapa} con coordenadas de 3 puntos turísticos
          \item Nombre del equipo y fecha
        \end{enumerate}
      \end{column}
      \begin{column}{0.42\textwidth}
        \begin{notabox}[Diseño]
          Recuerda:
          \begin{itemize}
            \item El número debe ser \textbf{preciso} y tener \textbf{unidades}
            \item El mapa debe tener \textbf{ejes} y \textbf{escala}
            \item El título debe mencionar \textbf{matemáticas}
          \end{itemize}
        \end{notabox}
        \dificultad{4}
      \end{column}
    \end{columns}
  \end{crearbox}
\end{frame}

% ── FRAME 13: Error frecuente ─────────────────────────────────────
\begin{frame}{¡Cuidado! «Más datos = mejor guía»}
  \begin{errorbox}
    \textbf{Error habitual:} incluir \textbf{todo} lo recopilado porque «más información es mejor».

    \vspace{0.3cm}
    \begin{columns}[t]
      \begin{column}{0.47\textwidth}
        \incorrecto\ \textbf{Guía con exceso:}
        \begin{itemize}
          \item 47 coordenadas sin clasificar
          \item Datos sin fuente mezclados con datos verificados
          \item Gráficas redundantes
          \item Números sin contexto
        \end{itemize}
      \end{column}
      \begin{column}{0.47\textwidth}
        \correcto\ \textbf{Guía bien seleccionada:}
        \begin{itemize}
          \item 10--15 datos \textbf{seleccionados y verificados}
          \item Cada dato tiene fuente y unidades
          \item Gráficas complementarias, no repetidas
          \item Números con contexto y comparación
        \end{itemize}
      \end{column}
    \end{columns}
    \vspace{0.2cm}
    \begin{center}
      \textit{Una buena guía tiene datos \textbf{seleccionados}, \textbf{verificados} y \textbf{relevantes}.}
    \end{center}
  \end{errorbox}
\end{frame}

% ── FRAME 14: Evidencia del día ───────────────────────────────────
\begin{frame}{Evidencia del día}
  \begin{evidenciabox}
    \textbf{Lista de datos seleccionados} (entregar al final):

    \vspace{0.3cm}
    La lista debe contener \textbf{exactamente 10 datos}, organizados por categoría, con:

    \begin{columns}[t]
      \begin{column}{0.55\textwidth}
        Para cada dato:
        \begin{itemize}
          \item El dato exacto (con número, unidades y fuente)
          \item La categoría de la guía a la que pertenece
          \item Puntuación en los 5 criterios R-E-C-R-O
          \item Justificación en una frase de por qué se incluye
        \end{itemize}
      \end{column}
      \begin{column}{0.42\textwidth}
        \begin{notabox}[Distribución mínima]
          \begin{itemize}
            \item 2 datos de Mapa
            \item 2 datos de Monumentos
            \item 2 datos de Rutas
            \item 2 datos de Gráficas
            \item 2 datos de Curiosidades
          \end{itemize}
        \end{notabox}
      \end{column}
    \end{columns}
  \end{evidenciabox}
\end{frame}

% ── FRAME 15: Resumen ─────────────────────────────────────────────
\begin{frame}{Resumen de la sesión 17}
  \begin{resumenbox}
    \begin{columns}[t]
      \begin{column}{0.48\textwidth}
        \begin{itemize}
          \item Hemos aplicado los \textbf{5 criterios} R-E-C-R-O para evaluar datos.
          \item Hemos aprendido a detectar \textbf{datos problemáticos}: sin fuente, sin unidades, obsoletos o implausibles.
          \item El \textbf{árbol de decisión} ayuda a tomar decisiones sistemáticas.
        \end{itemize}
      \end{column}
      \begin{column}{0.48\textwidth}
        \begin{itemize}
          \item La guía tiene \textbf{6 categorías} bien definidas.
          \item La \textbf{selección rigurosa} es parte del proceso matemático.
          \item En la próxima sesión: \textbf{construimos la guía completa}.
        \end{itemize}
        \vspace{0.3cm}
        \begin{formulabox}[BCrear]
          Calidad $>$ Cantidad
        \end{formulabox}
      \end{column}
    \end{columns}
  \end{resumenbox}
\end{frame}

\end{document}
